\documentclass[10pt]{article}
\usepackage[margin=20mm]{geometry}
\usepackage{amsmath,amssymb,mathtools,fontspec,unicode-math,xurl,hyperref}
\IfFontExistsTF{DejaVu Serif}{\setmainfont{DejaVu Serif}}{\setmainfont{Latin Modern Roman}}
\setmathfont{latinmodern-math.otf}
\setlength{\parindent}{0pt}\setlength{\parskip}{6pt}\sloppy
\title{Hölder Regularity and Smoothness of Complex Analytic Sets: Initial Technical Review}
\author{AI-assisted technical review}
\date{}
\begin{document}
\maketitle
Abstract: The regularity quantifiers, horn-neighborhood argument, projective tangent-cone step, separation of covering sheets, and sharpness example are examined. A sign error in the displayed log-Lipschitz modulus is established. The substantial implication from Hölder regularity for every exponent below one to smoothness does not use that erroneous expression. No Major mathematical defect is established in the inspected argument. The cited source evidence pages were visually checked in the recovery.\par
\section*{Bound source and coverage}
The work is José Edson Sampaio, Log-Lipschitz and Hölder regularity imply smoothness for complex analytic sets, arXiv:2404.06943v2, 29 October 2024. The corresponding publisher title adds the word regularity after log-Lipschitz; DOI 10.4171/JEMS/1721. The bound 23-page PDF has SHA256 prefix `6434ffcfac02e237`. Page numbers below refer to this PDF. The publisher’s 25 September 2025 publication date is not the preprint date.\par
All extracted pages and references were read. A distinct 26-page author-hosted PDF was consulted only at its introductory definition and section structure. It has different numbering and no established internal revision date; it is not substituted for v2. The arXiv record says v2 corrected a mistake in Proposition 3.2. That earlier correction is not itself a new finding.\par
The deeper checks concern Definitions 2.1–2.3, Theorems 2.14–2.15, Claims 3.1.1–3.1.8, Proposition 3.2, the diameter-distance proof, Theorem 5.1 and the scopes of Corollaries 5.2 and 5.5. External topology is traced at the steps below rather than represented as independently reproved in full.\par
\section*{Problem 1: the log-Lipschitz modulus has the wrong sign (Minor)}
\subsection*{Location and decisive contradiction}
On pp. 3 and 7 the displayed definition bounds \(\|h(x)-h(y)\|\) by \[C\|x-y\|\log\|x-y\|,\qquad C>0,\quad0<\|x-y\|<1/2.\] For every allowed distance \(\delta\) the right side is negative, whereas the left side is nonnegative, even for a constant map. For example, \(\delta=1/4\) gives \((C/4)\log(1/4)<0\). No map on a set having such a pair can satisfy this literal definition. A smooth complex line therefore fails condition (3) of Theorem 3.1 under that definition, despite satisfying (4) and (5). The separate author-hosted PDF retains the expression in its introductory passage; no corrected version was established.\par
\subsection*{Correction and impact}
Replace the modulus by \(C\delta|\log\delta|=C\delta\log(1/\delta)\). It is positive, and \(|\log\delta|\geq\log2\) on the specified interval, so every Lipschitz map satisfies the corrected bound after increasing the constant. For each fixed \(\alpha<1\), \[\sup_{0<\delta<1/2}\delta^{1-\alpha}|\log\delta|<\infty ,\] which proves the needed Hölder comparison.\par
The uncorrected definition and the literal equivalence component are false. The substantive proof of \((1)\Rightarrow(5)\) uses only Hölder charts and never uses the negative expression. The Minor assessment concerns a miswritten modulus for the identified regularity class; it is supported by this separation of dependencies, rather than merely by ease of editing. The evidence does not contradict the intended positive-modulus log-Lipschitz theorem or the all-exponent Hölder conclusion. No author adoption of the proposed correction is claimed.\par
\section*{Quantifiers and the tangent cone}
Definition 2.2 correctly distinguishes a possibly different chart for each \(\alpha<1\) from one chart valid for all \(\alpha<1\). The proof starts with the weaker hypothesis. The extension-and-shear construction on pp. 11–12 loses \(\alpha\) to \(\alpha^2\). Availability of exponents arbitrarily close to one absorbs this loss; neither a common chart nor a uniform Hölder constant is assumed.\par
In Claim 3.1.2, passage through a chart and back changes the horn exponent from \(b''\) to \(\alpha^4b''\). Fix \(b''>1\) in the stability interval, then take \(\alpha\) sufficiently close to one that \(1<\alpha^4b''<b''\). Changes in fixed horn amplitudes can be absorbed after restricting the germ. The step from bounding in a thicker horn to vanishing in the original horn uses the stabilized maps on homology, not just equality of group ranks.\par
This dependency was checked in the original research source Moderately discontinuous homology, arXiv:1910.12552v3, 22 January 2020, pp. 47–50: Theorem 112, the proof of Theorem 107, Remark 117 and Corollary 118. The direct limit is built from neighborhood maps, its stabilization uses homologically trivial cobordisms, and Remark 117 treats amplitudes. These passages support the naturality needed here. Their numbering differs from the published citation; editions are not conflated. The broader MD-homology theory and tangent-cone theorem remain external foundations.\par
For Claim 3.1.3, stabilization to \(N\geq k+4\) makes \(H^2(P,\widetilde V)\) and \(H^3(P,\widetilde V)\) vanish through the complement-homology bound. Thus the hyperplane class generates \(H^2(\widetilde V)\cong\mathbb Z\). The Gysin sequence propagates this generator to the top degree using the link cohomology vanishing. The fundamental-class pairing then forces degree one. Integral coefficients matter: rational rank alone would not suffice. Alexander duality, the affine-complement theorem and the Gysin sequence for the relevant singular base were traced, not fully reproved here.\par
\section*{Covering sheets and the sharpness example}
With a linear tangent cone, a proper projection is a finite covering away from its branch set. Two distinct lifts of a generic radial arc would have outer separation of order \(b>1\). Every connecting path must leave a branch-free conical sector, giving inner distance at least a positive multiple of \(t\). Connecting the two subanalytic lifts through zero gives the matching upper order \(t\).\par
A bi-\(\alpha\)-Hölder chart transfers a straight segment to a path of diameter at most a constant times \(\|u-v\|^{\alpha^2}\). The finite pancake argument gives \(d_{\rm inn}\leq mC\,d_{\rm diam}\). Hence two sheets would imply \(t\lesssim t^{\alpha^2b}\), impossible for \(\alpha^2b>1\). This checks the direction of the decisive inequality. Near-minimizing paths suffice in the noncompact version of the last-exit construction; exact geodesic attainment is unnecessary.\par
For Proposition 3.2 choose \(n>1\) and \(\alpha<b=n/(n+1)<1\) on \(x^n=y^{n+1}\). The link is a smooth embedded circle with two positive constant coordinate moduli, say \(a,c\). The weighted radial map \(F(tv)=(t f_1(v),t^b f_2(v))\) is \(b\)-Hölder by \(|s^b-t^b|\leq|s-t|^b\) and the Lipschitz link parametrization. For an independent inverse estimate, when \(0<t\leq s\leq1\) the first coordinate controls \(s-t\) by \(a^{-1}\) times the output distance, while \[t^{1-b}(s^b-t^b)\leq b(s-t)\leq s-t .\] Together these bound \(t\) times the difference of link points. The Lipschitz inverse link map then bounds the inverse radial map. Products with \(\mathbb C^{d-1}\) give the stated dimensions. No defect in this corrected v2 construction was established.\par
The statement of Proposition 4.1 displays the automatic direction \(d_{\rm diam}\leq mC\,d_{\rm inn}\), whereas its proof concludes the useful converse on p. 19. The statement should match that proved converse; the argument actually supplies the input to Theorem 4.2. The historical Theorem 1.1 should explicitly retain the complex-surface dimension from its preceding paragraph. These editorial observations are separated from the proved sign error.\par
\section*{Infinity, outcome and limitations}
Theorem 5.1 assumes algebraicity. Inversion takes asymptotic outer separation order \(b<1\) to \(2-b>1\), recovering the two-sheet contradiction. The pure-dimensional reduction and possible finite exceptional set are retained. Corollary 5.2 extends the result to entire analytic sets only for Lipschitz regularity, through a cited algebraicity theorem; it does not assert the all-Hölder statement for arbitrary entire analytic sets. Corollary 5.5 assumes a biholomorphism and uses holomorphic growth control for its affine conclusion.\par
One Minor definition error is recorded with exact repair and impact. No Major defect was established in the substantive inspected proof. This is not certification of every external theorem or every source page. A bounded publisher, arXiv and author-page search located no separate code or supplementary file. Known version corrections and the distinct author manuscript are retained as version information. No author code was run. The recovery inspected 18 source pages, the five identified MD-homology evidence pages and three introductory pages of the distinct author manuscript. Historical image-access failure is retained in the prior evidence. Final report-page viewing and compilation facts are recorded separately. Scientific decision is pending and human reviews are empty.\par
9 J. E. Sampaio, Log-Lipschitz and Hölder regularity imply smoothness for complex analytic sets, arXiv:2404.06943v2 (29 October 2024); journal DOI 10.4171/JEMS/1721. J. Fernández de Bobadilla, S. Heinze, M. Pe Pereira and J. E. Sampaio, Moderately discontinuous homology, arXiv:1910.12552v3 (22 January 2020), pp. 47–50. Published in Comm. Pure Appl. Math. 75 (2022), 2123–2200, DOI 10.1002/cpa.22013. J. E. Sampaio, publication page, https://sites.google.com/mat.ufc.br/edsonsampaio/research, consulted 27 October 2026; distinct author-hosted PDF retained in the assigned source manifest.\par
\section*{Supporting checks}
The original assigned assistant performed the textual review and independent sign, exponent and cusp arguments. The recovery executor reused the hash-bound report and checks, then supplied the missing source-image inspection and current report build and viewing evidence. No independent acceptance was performed. Controller-observed UI identity and runtime self-description are recorded separately, without asserting verified backend equivalence.\par
\end{document}
